\RequirePackage{fix-cm}
\documentclass[fontsize=11.5pt]{scrartcl}

\usepackage[letterpaper,hmargin=1.3in,vmargin=0.8in]{geometry}
\usepackage{microtype}
\usepackage{enumitem}
\usepackage{xcolor}
\usepackage[authoryear,round,sort]{natbib}
\usepackage[draft]{hyperref}
\usepackage{parskip}

\setlength{\parskip}{4pt plus 1pt}

% ---------------------------------------------------------------------------
% Usage, for each reviewer point:
%
%   \begin{reviewer}
%   Reviewer's comment, pasted verbatim.
%   \end{reviewer}
%
%   Our response.
%
%   \begin{revised}
%   Optional: new or changed text from the manuscript, quoted.
%   \end{revised}
%
% For reviewer text that needs no reply (sign-offs, reference lists), use
% \begin{reviewernote} ... \end{reviewernote}: same style, no number.
% ---------------------------------------------------------------------------

\definecolor{reviewercolor}{gray}{0.35}
\definecolor{revisedcolor}{rgb}{0.0,0.25,0.55}

% \reviewersection{Heading}{Label}: starts a reviewer's section; their points
% are then numbered Label.1, Label.2, ...
\newcounter{point}
\newcommand{\pointprefix}{}
\renewcommand{\thepoint}{\pointprefix.\arabic{point}}
\newcommand{\reviewersection}[2]{\section*{#1}\renewcommand{\pointprefix}{#2}\setcounter{point}{0}}

% Reviewer comment: numbered, indented, grey italic.
\newenvironment{reviewer}
  {\refstepcounter{point}\par\medskip
   \noindent\textbf{Comment \thepoint}\par
   \begingroup\leftskip=1.5em\rightskip=1.5em\itshape\color{reviewercolor}}
  {\par\endgroup\medskip\noindent\textbf{Response.}\ }

% Reviewer text with no reply needed: indented, grey italic, unnumbered.
\newenvironment{reviewernote}
  {\par\medskip\begingroup\leftskip=1.5em\rightskip=1.5em\itshape\color{reviewercolor}}
  {\par\endgroup\medskip}

% Quoted manuscript text: indented, blue.
\newenvironment{revised}
  {\par\begingroup\leftskip=1.5em\rightskip=1.5em\color{revisedcolor}}
  {\par\endgroup}

\begin{document}
\pagestyle{empty}

\begin{center}
{\Large\bfseries Response to reviewers}\par\medskip
\textit{Living Alignment}\\
PNAS manuscript 2026-22180
\end{center}

\vspace{1em}

% TODO: optional overview of the main changes.

% Reviewer comments below are imported verbatim from reviews.txt (typos
% preserved). The division of Reviewer 2's review into numbered points is ours.

\reviewersection{Editor}{E}

\begin{reviewer}
Both reviewers have valuable suggestions that must be taken seriously.
\end{reviewer}

We thank the editor and the reviewers for their careful reading and constructive comments. 

We found there was a great deal of overlap between the suggestions of the two reviewers. This central shared critique was highly useful for sharpening the message of the paper. 

To restate the central critique, the original paper was sometimes unclear about whether it was making descriptive, definitional, or normative statements. 


The problem surfaced in a few different ways:
\begin{itemize}
  \item Both reviewers directly noted this problem. R2 asked: ``Is this a phenomenological observation? Or is it an assertion about what you like (ie you dont like monoculture)? Is it a definition of what you mean [by] healthy, something correlated with healthy, or something that causes healthy?". R1 pointed out that the manuscript ``sometimes moves from how these properties support healthy functioning to these properties defining healthy functioning".
  \item Both reviewers were concerned about the is-ought gap. R1 noted the problem of distinguishing virtuous from vicious cycles, suggesting that generativity ``can't distinguish a flourishing system from one that generates harm in ever more elaborated and creative ways." R2 said, ``There isnt any objective argument to support the idea that a monoculture crop that `successfully' maximises yield is in some sense a failure".
  \item Both reviewers were concerned about a possibility of circularity. , and might get trapped into the ``idea that fixation is unhealthy and unhealthy means fixated". 
  \item Both reviewers were concerned about testability and falsifiability. R1 said it would be ``hard to see what predictions could fail from the authors' argument."
  \item R1 said that the paper touches ``on the importance of values [...] but not in a way that would allow that to be used effectively. In fact, the manuscript explicitly says that alignment is not about the right values or principles. That set up a kind of contradiction."

  \item 
\end{itemize}





Unresolved conflicts across levels. Treating health as a single unanchored property rather than situating it descriptively within nested systems meant we appeared to "refer to health at different levels of complexity without specifying how conflicts among these levels should be evaluated."


For Reviewer 2:

Unclear epistemic status. Reviewer 2 articulated the three-way ambiguity directly: "Is this a phenomenological observation? Or is it an assertion about what you like ...? Is it a definition of what you mean [by] healthy, something correlated with healthy, or something that causes healthy? The status of these observations [isn't] clear."

Circularity and the is–ought gap. Conflating a definitional claim with a normative one creates the danger of "we define healthy to be [X], so we should do [X]," raising the question: "Can we make it not a matter of opinion that this is a failure?" Otherwise, "one could be criticised for saying 'I don't like lack of diversity, so we shouldn't do things that kill diversity'"—since there is no objective argument that "a monoculture crop that 'successfully' maximises yield is in some sense a failure. It is a success on its own terms."

Unfalsifiability. Mirroring Reviewer 1's concern about predictions, drifting into definitional or normative territory meant that "There [aren't] a lot of statements here that are falsifiable."

Attitudinal opinion rather than mechanism (and reducing to Goodhart's law). Whenever the paper read as a normative exhortation rather than a mechanistic synthesis, it appeared "supported loosely by some anecdotal narrative," where "there is no mechanism proposed to do it. Its recommendation is 'attitudinal' - a direction of travel" ("How do we do that, what is the mechanism?"), and risked a reading in which the message simply "reduces to 'all things in moderation (except moderation)'" or "basically advocates for Goodhart's law."



In addition to this shared core issue, the reviewers raised other valuable points, like the role of contingencies in reward hacking, connections to multilevel selection and Ostrom's design principles, and additional important citations. We address these in the point-by-point below.



% \begin{itemize}
%   \item Descriptive / mechanistic. There are commonalities across systems in terms of the properties that are widely associated with health.
%   \item Definitional. Avoiding fixation is what health and alignment mean. For example, we said things like, ``we've defined aligned systems by...".
%   \item Normative: Ongoing contextualization (with increasing subtlety and sensitivity) is inherently ``good", and we ``should" align to it.
% \end{itemize}

% Without anchoring health to an independent functional outcome, fixation appeared to be a ``just a topographical" concept rather than a ``functional" one. 

\reviewersection{Reviewer 1}{R1}

\begin{reviewernote}
There is a lot to like about this ``perspectives'' paper. It argues persuasively that any fixed formalism in the areas of alignment can go wrong if relied upon too rigidly. The core move - flexibility at every scale - makes sense and fits with a large body of literature. And the authors should also be given credit for raising some of the hardest questions, including ``healthy for whom?'' It will foster needed discussion in the area of AI alignment.

There are some weaknesses in scope and substance, however, which I would like to share more as possible ways the authors might want to strengthen their piece. My biggest concern is that the evolutionary thinking coming through the paper stops half-way.

For example, they cite Kashdan and Rottenberg (2010) on psychological flexibility but keep only the ``moving between frames'' part of the larger argument. The authors touch on the importance of values, which is empirically key to the positive impact of psychological flexibility (Hayes, 2019), but not in a way that would allow that to be used effectively. In fact, the manuscript explicitly says that alignment is not about the right values or principles. That set up a kind of contradiction.

You can see the more general impact of the ``half step'' issue in how the paper addresses health, which is defined almost entirely by what it avoids: fixation and collapse. As an evolutionary matter that is like having variation and replication/retention without a selection criterion. If you skip the hard work of addressing selection (hard work because ``survival'' alone simply will not work once we are applying evolutionary principles to human affairs, as this paper explicitly is attempting to do), ``generativity'' can't distinguish a flourishing system from one that generates harm in ever more elaborated and creative ways. A vicious cycle is evolution gone awry; a virtuous cycle is something else entirely. Generative complexity could feed either. That is a more technical way of saying exactly what is worrying people about AI. The paper as written can't distinguish the two types of cycles.

Without a selection criterion some parts of the authors' argument comes precariously close to circular reasoning. The idea that fixation is unhealthy and unhealthy means fixated would be such an argument. Similarly, the manuscript makes a good cse that semipermeable barriers, modularity, and flexibility can contribute to the adaptive function of living systems. But it sometimes moves from how these properties support healthy functioning to these properties defining healthy functioning.

There is a ready way to break out of circles such as these: specify the selection criterion. Then, for example, narrowing becomes fixation only when it stops working for the specified purpose or value. Stated another way, fixation then becomes a functional concept, not just a topographical one, which is far more in tune with the broad themes of this article.

The author's arguments could use an organizing framework and one is readily at hand. The authors acknowledge that what is healthy for a cell may not be healthy for a multicellular organism and what benefits and organism might harm an ecosystem. They ask whether AI could destroy humanity while creating something extraordinarily generative. They refer to health at different levels of complexity without specifying how conflicts among these levels should be evaluated.

Multilevel selection theory provides a natural framework in such areas. MLS research has shown that a process can be adaptive at one level of complexity and destructive than another. Suppressing destructive competition within groups facilitates cooperation, but cooperation within groups can also intensify destructive competition between groups. Elinor Ostrom won a Nobel Prize for showing how core design principles can help manage these kinds of conflicts - and do so in ways that are very much in touch with the author's arguments (e.g., Wilson, Ostrom and Cox, 2013).

Moving in this direction would not take that much space and it would have other positive effects on the manuscript. The key points in the paper are often presented only by analogy. The paper succeeds anyway because the examples are very well chosen, but it will limit the long-term impact of the article if it is left at that level. For one thing, it's hard to see what predictions could fail from the authors' argument \ldots\ and researchers sensing that will likely not test the ideas this piece contains. That would be unfortunate. A multi-level selection approach would make ready sense of almost all of the author's best empirical examples and do so in a way that is more testable and less metaphorical and analogical.

A brief comment in another area: reward hacking. From a contingency standpoint, a system that games its reward is just doing what its contingencies selct for, like a child learning to lie. Hortatory rules backed up by aversive contingencies tend to teach concealment, not honesty. Authoritarian parents learn that the hard way, and I fear AI researchers are learning it anew. I appreciate that the authors are going beyond that with their comments on precommitment, cognitive control, internal regulation and the like but more of a focus on purpose, values, and the contingencies that establish and maintain these processes could carry the broader argument additional steps forward.

I think this is a strong article, and I hope the comments above strengthen it further.
\end{reviewernote}

\noindent\textbf{Response.} We were happy to read Reviewer 1's overall agreement with our viewpoint on flexibility across scales and the importance of this message in the current AI landscape. We are grateful to them for raising their thoughtful and pertinent concerns.

First, we want to try to restate the Reviewer's comments, to make sure we've understood them well. It seems the central concern is that our framework could use a clearer external anchor to evaluate what is healthy. For example, they suggest that narrowing becomes fixation when it stops working for a given purpose or value. They see generativity, flexibility and complexity as being able to go either in a positive or negative direction (a virtuous or vicious cycle). The Reviewer points out that for Kashdan and Rottenberg (2010) and for Hayes (2019), psychological flexibility is primarily an instrumental means toward an end of values or purpose -- not an end in itself. Purpose provides an anchor so flexibility can be useful `for' something. Relatedly, the Reviewer is concerned about a risk of circular definition, about a possible contradiction between the importance of values and our claim that alignment is not about having the right values, and that the framework might not be testable. 

The Reviewer suggests that a solution would be to specify a selection criterion. They also observe that multilevel selection (MLS) addresses conflicting pressures between levels; because our framework is about contextualization (how entities play a partial role in larger structures), MLS would fit naturally. In particular, the Reviewer notes Ostrom's design principles as a way of managing conflicts between levels. Finally, the Reviewer notes that reward hacking, which the paper uses as a central example, is a direct result of contingencies: if the agent is rewarded for giving the right answer, but the easiest way to give the right answer is to cheat, then it's not surprising that the system learns to cheat. They invite us to discuss in more depth how contingencies can be configured to establish and maintain longer-sighted mechanisms like precommitment. 

We hope that we've accurately interpreted the Reviewer's points. 


*** Below here is just early vomit draft notes, we're not committed to any of it ***

Option 1: fully punt on values. We would double down on the idea that we're simply identifying commonalities in the properties of systems that are considered healthy. For example, it's widely understood that depression or Parkinson's disease is 'unhealthy'. In this case, we wouldn't be trying to define health or trying to define what is good. We would simply be reporting descriptively on the properties of systems that are widely understood as being 'healthy'. Then, in the alignment section, we wouldn't make any value judgement about alignment: we would simply say that to whatever extent someone wants systems to be healthy, they could follow the principles we've described. This is purely an 'is' argument, in Hume's language, rather than an 'ought' argument.

Option 2: take a stance on values. This is trickier, because a central part of our argument is that any particular well-defined value system is misaligned. But we would have to take a position that there is something 'good' that intrinsically cannot be well-defined. (For this we could lean on many philosophers, including Pirsig.) In our language, what is 'good' is something about continually increasing subtlety, releasing into larger identification, becoming sensitive to more and more structure. In this second option, we would run a significant risk that people would say, 'you're making your argument totally nonfalsifiable by saying that you can't define the thing you're talking about'. This position is hard to defend. On the other hand, it could also be rich and powerful.

In response to the Reviewer's central concern -- that our framework could use a clearer external anchor to evaluate what is healthy -- we have made the following improvements. 

* We make it more clear that the contribution of Section 1 is to link existing notions of `health and flourishing' from across fields, using a common language. In this section, we're not trying to set up our own value system about what health is. For example, it's commonly accepted that depression, Parkinson's disease, a cell's death or a civilization's collapse are `unhealthy' for the systems suffering from them. In some cases the examples are looser, like claiming that a society is healthier if it's continuing to solve its own problems and make new discoveries. But generally we're looking at relatively uncontroversial pre-existing examples of `healthily-functioning' or `flourishing' systems, taken broadly. 

They all look like instances where one mode of interaction gains more traction. Healthy systems thus consist of parallel threads exploring different modes of interaction across time. And they inevitably include barriers to limit modes of interaction. This is just a descriptive claim about a simple way of describing commonalities across systems. 

* General direction of increasing nuance.

* link to alignment

OR, do we want to say, even in Section 1, that we do have this value system, that healthy life is this openness to increasingly subtlety? Or, we could start with the descriptive thing and then add the value thing?

OR, another thing we could do is separate out the straightforward examples (like addiction, obsession). Give those first. Examples where health looks like avoiding collapse or fixation. Then, argue for a broader principle (getting into everything R1 and R2 talked about).




Some of the points we fully agree with and have incorporated into the paper. Other points have given us inspiration to make our distinct arguments points more clearly. We divide our reply into these two categories below.


Committing to a particular anchor would be fixating. 

The `good' direction is increasing sensitivity/subtlety/long-sightedness. Not the same as complexity or flexibility. By construction it doesn't have a single formalization. The idea is that whatever 

``Clearly defined boundaries. All examples of major evolutionary transitions involve groups with clear boundaries, such as the cell wall for cells and nests for eusocial insects." -- from the Wilson, Ostrom and Cox (2013) paper


First, we agree that flexibility is often instrumentally useful toward particular ends like certain notions of well-being. 

What we're proposing is that these life-like properties are actually what we mean, aross different kind of systems, when we talk about health. So, we're not trying to use health to define flexibility, but we are trying to use flexibility to operationalize what we mean by `health'. For example, . We actually mean that they retain this potential. So then it comes to the question of, `what about unhealthy forms of complexity'? We suggest that, for any given system, what defines its health is that it maintains the kind of rich, flexible potential that we have described. Of course, smallpox's health might be a human's fixation, and vice versa. 


We fully agree about the importance of values. We also stand by our claim that alignment is not about the right values or principles. We don't believe this is a contradiction. Not committing to a fixed formulation of values doesn't imply that values have no guiding role. In fact, our perspective is that there is *nothing but* local value systems to guide us.

In terms of reward hacking and contingencies, the paper attempts to bridge or generalize between evolutionary time scales and learning time scales. Precommitment, cognitive control, and so on are examples of barriers that cause some more myopic process to be contextualized within some larger process. In general, what are the contingencies or mechanisms that could establish beneficial barriers like precommitment. The general answer looks something like multilevel selection -- because it's like the pressure from the outer (or contextualized, or larger) system, which is limiting the myopia of the smaller system. This is the mechanism (or contingency) that brings about barriers. The reasons we don't say it's exactly multilevel selection: 1) there are mechanisms other than selection, like a human creating a painting, 2) there's no final selection criterion

We don't think there is an absolute, universal anchor. But there is an interesting sense in which we can look at nested inclusivity as a step beyond pure relativism. This is why it would likely be misaligned, in our proposal, for AI to wipe out humanity. 

What we're proposing isn't that any particular notion of flexibility would be an end to itself. In fact, any particular kind of flexibility would be unhealthy under our definitions, if fully committed to. 


We're also proposing that alignment not be ultimately anchored in any particular notion of harm versus virtue. 


Because the Reviewer's constructive points are interrelated, we decided to try to address them holistically.


The Ostrom-style barriers are good examples. She's mainly looking at social systems. Ostrom's design principles are excellent examples of barriers against fixation. 

I don't think our thing is really about selection. For a couple of reasons. One, selection doesn't necessarily imply flourishing. In the end, the heat death of the universe is selected for. Or paper clips could be selected for. And two, I don't think it really has to be selection in the normal sense of selection out of some pool of variation. It could be some other method of de novo construction. Like, artificial selection would be if there were 100 cows and the humans chose one of them that they liked. But if I just draw a painting, I haven't really "selected" that painting.

For any definition of diversity, if you maximize that, it's also fixation.



Our paper may seem to be taking a position that flexibility is an end to itself.

We would like to disentangle two ideas. We agree that *having* values is key to the impact of flexibility. Flexibility is often locally in service of something. But we contend that this doesn't make any particular set of values right or wrong in an absolute sense. Within one value-frame, flexibility can be beneficial. But even a broad value-frame (like some conceptualization of 'well-being') is still a frame. There's a more subtle (more richly contextualized) frame from which excess attachment to that conceptualization is a lack of flexibility. We the authors find that we tend to value this process of stepping back into more inclusive frames, but we also acknowledge that not everyone values this.

Our perspective is that the notion of 'well-being' itself is not something that can be fully pinned down.

Not a fixed notion (like a definition of well-being) to align to,

This is why the open-ended perspective is so important as AI is optimizing the hell out of objectives and bringing about faster and faster changes. In the past, we could rely on approximate intuitive notions of well-being. But now we realize that pinning down *any* definition of well-being would lead to fixated outcomes.

We intended for the paper to make exactly this distinction: between flourishing and ever more elaborated harm. It's an unusual argument, so we have taken the opportunity of this revision to try to make it much more clearly.

There's one value system where harm is in fact exactly fixation. This is quite a powerful and underappreciated one, which can unify some of the kinds of values that people tend to deeply believe in.

But it's still just a conceptual lens.


In terms of the question `good for whom?', one nice example is assortative mating. Preferring similar mates is a barrier that causes a species to diversify from others. Interestingly, it creates a type of local collapse. You could think of this as exploitation, but it's probably more like exploration moved inward.  It's not all bad, probably essential to the living process.  That myopia can find new things.  However, if it is the only myopic process going on on earth, we'd say this approach is a bit risky.  It's probably good to have diverse context outside that system (same level or greater).  This doesn't just create parallel paths for selection to pick from, but also creates latent potential for combinatorial emergence (sex, introgression, horizontal gene transfer, endosym).


First, without getting into values, we do want to assemble evidence that what people generally mean by `unhealth', across many kinds of systems, can be functionally understood as fixation.

So the first point is that fixation is a powerful tool for understanding what `unhealth' \textit{is} at a more mechanistic level.

The second point is that this can be an interesting way to think about alignment. An open-ended way of approaching alignment, without being forced into eternally committing to a value scheme. You could think of it as `bypassing' values (after all, `values' is on some level just a human concept). Or you could think of it as more fully honoring all existing and potential future values.

The question about the relevance of clear values and selection criteria is helpful and nuanced. Our view about the fundamental importance of avoiding rigid adherence to values revolves around a few simple principles:
\begin{itemize}
  \item People can be fairly well approximated as having value systems.
  \item For any given value system, if it's taken too seriously, the world becomes locally less life-like.
  \item By definition, a less life-like environment is less likely to support the flourishing of life as we know it.
  \item We acknowledge that we the authors have value systems, conscious and unconscious: We may be attached to something about life in its current form, we may want recognition and to contribute something meaningful, we may be excited about something in these ideas and their magnetic pull on our psychology, etc
  \item In spite of this, we urge caution about claiming that lifelikeness or any other specific values are the right alignment target in some absolute sense.
\end{itemize}

For this reason, we feel that being too prescriptive about particular values in the paper may interfere with our overall message and deviate us away from the viewpoint we are trying to advocate for.

In the revision, we've taken on board the reviewer's advice about adjusting the text more in line with Multilevel Selection Theory. This theme was discussed by the authors during writing but we found the topic so deep and complex that it was ultimately excluded from the submitted draft. The reviewer's comments have made us reconsider this decision, and we now include and integrate MLS into the broader theoretical landscape in sections [X, Y, Z].

It's more lifelike when it's more contextualized; more inclusive. In evolution, the fact that selection pressures operate at a larger scale may partly explain this (eg selection for a whole group)

This can explain how things get more longsighted (in time and space). More longsighted in a general topological sense means less fixated.

However, we are neutral on the question of whether this is the right framework to think about psychology and society. Some authors take the stance that because ideas and cultures and institutions spread and succeed or fail, we can use survival and selection logic to understand them.

We suggest that multilevel selection is a special case of

Living systems have many individual parts that have, over long periods of time, faced situations from many, many different angles. This creates a rich repertoire. This kind of repertoire tends to generalize better to future and long-term and out of distribution situations.

For example, people deliberately invent strategies that contextualize their own impulsive tendencies. Societies contextualize the local drives of individuals. These don't have to be selected for in a survival sense.

MLS is related to the point about contextualization. Every local perspective is contextualized. MLS conveys longer-sighted optimization.

Some versions of MLS do apparently also account for heterarchy (where entities are nested in multiple kinds of systems so there aren't levels).

Short-sighted optimization leads to long-term trouble \citep{dawkins1976selfish, traulsen2006evolution, wilson1980natural, wagner1996perspective}. For example, pathogens perform better in the short-term if they voraciously consume the host. But if they kill their hosts too quickly, the infection won't spread in the population, and the pathogens will eventually perish. Bacteria have evolved mechanisms to recognize and suppress defectors among themselves \citep{west2006social}. These mechanisms are boundaries that contextualize the myopic interests of individual pathogens, allowing the pathogen community as a whole to prosper in the longer term. Cancer follows a similar pattern. Individual cells in an organism can maximize the short-term success of their own lineage by replicating without respect for the tissue or the organism -- often leading to the demise of the entire organism. The organism has boundaries to prevent cancers from starting, such as the intracellular p53 system that detects DNA damage and can initiate apoptosis; and mechanisms to destroy cells from the outside when they overrun their own internal defense systems \citep{hanahan2000hallmarks, schreiber2011cancer}. These boundaries ensure that the whole organism can reproduce in an organized way, with longer-term benefits.

I think there are three things in what R1 is saying overall: one is about values, one is about multi-level selection, and one is a broader question about the mechanisms that give rise to non-fixated systems.

We agree that reward hacking is a canonical example.

I think the reviewer is asking, what are the *mechanisms* by which precommitment, cognitive control, and internal regulation come about.

At the level of abstraction of the paper, barriers are the proposed mechanism. But we can go one level farther and say that there is some kind of `external pressure'. We think the various multi-level selection models are mostly special cases of this template.

We are deeply thankful to Reviewer 1 for their inputs---their contributions have significantly improved our paper. The inclusion of the MLS work has helped to bolster our central message while also clarifying how varied selection criteria operate across multiple levels of life to create the complex life-like balances that pervade the natural world.

% ===========================================================================
% Ratings: Suitable quality: No. Sufficient general interest: Yes.
% Conclusions justified: No. Clearly written: Yes.
\reviewersection{Reviewer 2}{R2}

\begin{reviewer}
The paper advocates for the position that life does not optimise for any simple objective, and that doing so is antithetical to the creative processes of living systems. From this, the suggestion is that AI, if we want it to make a positive contribution to creative fourishing, also should not be directed at goal-achievement/objective maximisation either. In brief, any one criterion of alignment is misaligned.
\end{reviewer}

[Response.]

\begin{reviewer}
I feel torn in writing this review.

On the positive, I completely agree with the message of the paper. Its correct, important, not widely understood and its vital that more people undertand it. On the negative, I struggle to see it as a PNAS paper - its an opionion, supported loosely by some anecdotal narrative. There is `finding' or `discovery', there are no methods used, and there is no mechanism proposed to do it. Its recommendation is `attitudinal' - a direction of travel. I think its important, but I can imagine it being a piece in Noema rather than a PNAS paper. Perhaps, PNAS could take the position that, although it doesnt usually publish pieces that dont have methods or describe mechanisms, finding or discoveries, perhaps it should. Indeed, part of the problem this paper talks about is a kind of fixation on `whats it good for?'/`what can I do with that?'/`what have you proved?' that makes it more difficult to\ldots\ stop and smell the roses, ask `why are we doing this?'/`whats really important here?'. So, not publishing it is, in a sense, consistent with the problem that the paper identifies.
\end{reviewer}

[Response.]

\begin{reviewer}
There is a bit of a danger that the message of the paper reduces to `all things in moderation (except moderation)' which, although wise, is hardly new - It basically advocates for Goodhart's law - and the paper doesn't offer concrete ways to operationalise that (to alleviate the greed that drives the Western worldview/market-based infrastructure). Again, I'm totally on-board - but unfortunately, the problem isn't just one of how to align AI, we have to deconstruct to entire Western narrative. We live lives that are goal-directed, aligned with fixed objectives, and generally don't really that's a mistake until we get to middle-age - and its a societal problem not just a personal problem (or an AI lignment problem).
\end{reviewer}

[Response.]

% Note: the first sentence below has gaps as received ("healthy to be ,"
% and "should do ."). The reviewer's placeholders, probably in angle
% brackets, seem to have been stripped out by the submission system.
\begin{reviewer}
There is a bit of a danger that the paper becomes: we define healthy to be , so we should do . Where in this case is not a monoculture that has lost diversity, for example. The examples show that single-minded objectives lose creative diversity\ldots\ but they do meet their objectives. Can we make it not a matter of opinion that this is a failure? I mean, as is, one could be criticised for saying `I don't like lack of diversity, so we shouldn't do things that kill diversity'. I agree. But can it be converted into any sense that killing diversity is an objective failure, so you make an argument for the contrary instead of just stating a purpose? Im not sure it can, to be honest. It is just a preference. There isnt any objective argument to support the idea that a monoculture crop that `successfully' maximises yield is in some sense a failure. It is a success on its own terms. And that is the problem. Goodhart's law.
\end{reviewer}

It succeeds locally (on its own terms), but it's a failure in larger terms. 

\begin{reviewer}
I agree with most of the observations made about example systems.

There arent a lot of statements here that are falsifiable.
\end{reviewer}

[Response.]

\begin{reviewer}
You might mention the permaculture literature in the ecosystem section (I didnt follow your refs, maybe they are).
\end{reviewer}

[Response.]

\begin{reviewer}
Group problem solving in evolution has some recedent in hybrid vigor and shifting balance (for what it's worth).
\end{reviewer}

[Response.]

\begin{reviewer}
On creativity through compositional discoveries you might mention Boden.
\end{reviewer}

[Response.]

\begin{reviewer}
I love that bit about `for an agent built inside a frame the frame is invisible' (curious you didnt mention Godel).
\end{reviewer}

[Response.]

\begin{reviewer}
At the end of that section Im thinking about `holding a pov' (in a manner sustained against evidence) and `holding contradictions' (allowing opposites to be both true). Interesting territory.
\end{reviewer}

[Response.]

\begin{reviewer}
Mention MArgulis for endosymbiosis!
\end{reviewer}

[Response.]

\begin{reviewer}
Some of this sounds like straightforward explore vs exploit.
\end{reviewer}

[Response.]

\begin{reviewer}
A question that is bubbling up now: is moderation algorithmic? or non-algorithmic?
\end{reviewer}

[Response.]

\begin{reviewer}
Top of page 7, `healthy living systems have multiple semi-autonomous perspectives loosely and modularly coupled at many scales' - I agree. Is this a phenomenological observation? Or is it an assertion about what you like (ie you dont like monoculture)? Is it a definition of what you mean my healthy, something correlated with healthy, or something that causes healthy? The status of these observations isnt clear. I think its correlated with healthy (not a definition) and I dont think its causal - like, I cant take that as a recipe I can implement and get healthy life as an output. Id like to, its a good research goal, but stating as an aseertion/observation doesnt help me do that.
\end{reviewer}

[Response.]

\begin{reviewer}
I like the bit about `stretching below language'.

The decentering concept also very useful - thanks for that!
\end{reviewer}

[Response.]

\begin{reviewer}
`the continued dance of contextualizing and particular form' - beautiful. Compeltely agree. How do we do that, what is the mechanism?
\end{reviewer}

[Response.]

\begin{reviewernote}
I hope some of that is useful.
\end{reviewernote}

\bibliographystyle{abbrvnat}
\bibliography{../../assets/proxyfailure}

\end{document}
